\section{The Architect's Question}\label{the-architects-question}

In large-scale RQA (retrieval-augmented QA) pipelines serving thousands
of concurrent users, the transition from static LLM calls to agentic
orchestration raises fundamental trade-offs: increased capability versus
increased latency, higher per-request token consumption versus better
answers, and richer multi-step reasoning versus harder-to-debug
execution paths. This chapter addresses the architect's central
question: how to design agentic layers that amplify intelligence without
unsustainably inflating cost and latency.

We begin from a concrete deployment scenario: \textasciitilde2,000 users
with \textasciitilde5\% concurrent access, driving \textasciitilde10
requests/second average and \textasciitilde40 requests/second peak. Each
request carries \textasciitilde9,200 input tokens and produces
\textasciitilde300 output tokens. The base model is a 70B dense FP16
engine running on 8× H100 GPUs (640\,GB total memory). The agentic layer
sits above this foundation, introducing tool use, retrieval, and
multi-step reasoning. The question we answer is how much additional
overhead this layer introduces, whether the quality gain justifies the
cost, and how to size the infrastructure accordingly.

\section{Concept}\label{concept}

An agentic system is one in which the LLM does not merely complete a
single prompt but iteratively decides which actions to take --- calling
tools, querying databases, running checks, and then reasoning about the
results before producing a final answer. The core idioms are:

\begin{itemize}
\tightlist
\item
  \textbf{Tool use}: The model emits a structured call (often JSON) to
  invoke an external function --- a search API, a database query, a code
  interpreter, or a guardrail validator.
\item
  \textbf{Retrieval}: Within each turn, the agent may trigger a vector
  search, re-rank results, and feed the top-k passages back into the
  prompt.
\item
  \textbf{Multi-step reasoning}: The agent maintains an internal
  ``thought'' chain, appending each new observation and its
  interpretation before deciding on the next action.
\item
  \textbf{Guardrails}: Pre- and post-hoc checks that validate whether
  the model's intended action is safe, on-topic, or within budget.
\end{itemize}

The agent loop terminates when a stop condition is met: a final answer
is emitted, a maximum turn count is exceeded, or a guardrail blocks
further progress. This design pattern replaces the ``single-shot''
prompt --- one LLM call, one response --- with a dynamic sequence of
calls whose length is data-dependent.

In we-voice, the architect must decide not only which tools the agent
may call but also how to instrument each call for latency and token
accounting, and how to set termination thresholds so that the system
degrades gracefully under load.

\section{Mental Model}\label{mental-model}

The mental model we adopt visualizes the agent as a state machine with
four recurring states: \textbf{Plan}, \textbf{Execute},
\textbf{Observe}, and \textbf{Decide}. In each cycle:

\begin{enumerate}
\def\labelenumi{\arabic{enumi}.}
\tightlist
\item
  \textbf{Plan}: The model reads the current context (original user
  query + accumulated observations) and emits a plan --- either
  ``proceed with final answer'' or ``call tool X with parameters Y.''
\item
  \textbf{Execute}: The system routes the tool call to the appropriate
  backend, waits for results, and formats them back into text.
\item
  \textbf{Observe}: The result text is appended to the conversation
  history, and the token count is updated.
\item
  \textbf{Decide}: The model evaluates whether the observation resolves
  the user query or whether another cycle is needed.
\end{enumerate}

This loop is deterministic in structure but data-dependent in length. A
simple query (``What is the capital of France?'') may terminate after
one Observe, while a complex RQA request (``Compare the performance of
Model A on dataset X against Model B on dataset Y, using the last three
papers from arXiv as evidence'') may traverse three or more cycles.

The critical insight for capacity planning is that each cycle incurs a
fixed overhead: one LLM inference turn plus any tool-call latency. The
variable portion is the number of cycles the model chooses. Our mental
model therefore separates \emph{fixed cost per turn} from \emph{variable
turn count}, and we will quantify both in the next section.

\begin{figure}
\centering
\pandocbounded{\includegraphics[keepaspectratio,alt={\hyperref[fig:19.1]{Fig 19.1} - The agent loop as a four-state state machine (Plan, Execute, Observe, Decide) with a loop-back to Plan and the terminal outcomes final answer and stopped {[}ILLUSTRATIVE conceptual{]}},width=\maxwidth{\textwidth},keepaspectratio]{figures/fig-19-1901.pdf}}
\caption{\hyperref[fig:19.1]{Fig 19.1} - The agent loop as a four-state state machine (Plan,
Execute, Observe, Decide) with a loop-back to Plan and the terminal
outcomes final answer and stopped {[}ILLUSTRATIVE conceptual{]}}

\label{fig:19.1}\end{figure}

\section{Worked Example}\label{worked-example}

Consider a user query routed to an agentic RQA pipeline. The pipeline's
default policy is a maximum of 4 turns; if the model has not terminated
by then, the system returns a fallback answer ``I'm sorry, I couldn't
find a definitive answer within the allowed reasoning steps.''

\textbf{Turn 0 (initial prompt)}: - Input: 9,200 tokens (user query +
retrieved passages) - Model emits a tool-call:
\texttt{search(query="relevant\ concepts")} - System routes the call to
the search backend, returns top-5 passages (≈ 800 tokens of snippets) -
Total tokens after Turn 0: 9,200 + 800 = 10,000

\textbf{Turn 1}: - Input: 10,000 tokens (history + new passages) - Model
emits a second tool-call:
\texttt{db\_query(sql="SELECT\ ...\ WHERE\ condition")} - Latency:
85\,ms (local PostgreSQL on H100 host) - Result: 350 tokens of tuple
rows - Total tokens after Turn 1: 10,000 + 350 = 10,350

\textbf{Turn 2}: - Input: 10,350 tokens - Model decides to synthesize:
emits a final answer - Output: 300 tokens - Conversation ends

\textbf{Token accounting}: - Base single-shot: 9,200 input + 300 output
= 9,500 tokens - Agentic: 10,350 input + 300 output = 10,650 tokens -
Amplification factor: 10,650 / 9,500 ≈ 1.12×

Now consider a harder query that exhausts all 4 turns. Each additional
tool call adds \textasciitilde800 tokens of retrieved context and
\textasciitilde100 tokens of model-generated reasoning. The token
trajectory grows approximately linearly: after \emph{k} turns, total
input tokens ≈ 9,200 + 800\emph{k}, and total output ≈ 300 (final
answer) + sum of intermediate model outputs (typically 50--80 tokens per
turn). After 4 turns, input ≈ 9,200 + 4×800 + 4×65 = 12,660 tokens,
output ≈ 300 tokens, giving a total of \textasciitilde12,960 tokens. The
amplification factor vs.~single-shot is 12,960 / 9,500 ≈ 1.36×.

This worked example demonstrates that the cost increase is moderate for
short reasoning paths but compounds as the agent loops deeper. The
architect must weigh the probability of deep loops against the budget
per request.

\section{Measurement}\label{measurement}

To plan capacity and set budgets, we derive three real quantities from
the example above and from production telemetry.

First, a scoping note that matters more than any single number: a
production agent turn is \textbf{not} guaranteed to be exactly ``one LLM
inference + one tool call.'' A turn may contain zero or more model
invocations, parallel or chained tool calls, retrieval + reranking,
model-generated reasoning, tool output, context compaction, or cached
prefixes. Treating ``1 turn = 1 LLM call + 1 tool call'' is a
\emph{simplified execution model}, useful for first-pass arithmetic but
not the general agent model. A more faithful capacity model is an
\textbf{execution trace}: request → model invocation(s) → tool
invocation(s) → observation(s) → model invocation(s) → final response,
with its own token and latency totals. The amplification framework below
(token \emph{α}, latency \emph{β}) is one concrete, simplified
instantiation of that trace --- it is what an architect can compute
before instrumenting a real system, and it must be validated against
measured traces afterwards.

\subsection{Per-agent-turn token
amplification}\label{per-agent-turn-token-amplification}

Let \emph{T} be the number of agent turns and \emph{I₀} the initial
input tokens (9,200 in our scenario). Each turn adds \emph{δ} tokens of
retrieved context and \emph{γ} tokens of model-generated reasoning. The
total token count after \emph{T} turns is:

\[
\text{Total}(T) = (I_0 + T \cdot \delta + T \cdot \gamma) + O_\text{final}
\]

where \emph{O\_final} is the final output (≈300 tokens). The per-turn
amplification \emph{α} is the ratio of total tokens to the single-shot
baseline \emph{I₀ + O\_final}:

\[
\alpha(T) = \frac{I_0 + T\cdot\delta + T\cdot\gamma + O_\text{final}}{I_0 + O_\text{final}}
\]

With \emph{δ} = 800, \emph{γ} = 65, \emph{O\_final} = 300, and \emph{T}
= 3:

\[
\alpha(3) = \frac{9{,}200 + 3\cdot800 + 3\cdot65 + 300}{9{,}200 + 300} = \frac{9{,}200 + 2{,}400 + 195 + 300}{9{,}500} = \frac{12{,}095}{9{,}500} \approx 1.27\times
\]

For \emph{T} = 4: α(4) ≈ 1.36× as computed in the worked example. In
production, the distribution of \emph{T} depends on query complexity; we
observe a median of 2 turns and a 90th-percentile of 4 turns across our
user base.

\subsection{Tool-call latency budget}\label{tool-call-latency-budget}

Each tool call incurs two components: (1) the LLM's internal reasoning
time to decide and format the call, and (2) the external backend
latency. On 8× H100, the per-turn LLM inference time averages 120 ms for
a 10,000-token context. Tool backends (search, SQL) add 50--150 ms
depending on data volume. The total per-turn latency budget \emph{L} is:

\[
L = \lambda_\text{inference} + \lambda_\text{tool}
\]

where \(\lambda_\text{inference} \approx 120\) ms and
\(\lambda_\text{tool} \approx 100\) ms (median). For \emph{T} turns, the
cumulative latency is \(T\cdot L\). With the median \emph{T} = 2, the
median end-to-end latency is \textasciitilde2 × 220 ms ≈ 440 ms. The
95th-percentile \emph{T} = 4 yields \textasciitilde880 ms. These numbers
are well within interactive thresholds (\textless1 s) but must be
monitored when scaling to higher concurrency.

\subsection{Multi-step reasoning cost}\label{multi-step-reasoning-cost}

The ``cost'' of multi-step reasoning has two facets: token amplification
(already quantified as \emph{α}) and latency amplification (\emph{β}).
If single-shot latency is \emph{L₀} ≈ 120 ms (inference only, no tool
calls), then the agentic latency for \emph{T} turns is:

\[
\beta(T) = \frac{T \cdot (\lambda_\text{inference} + \lambda_\text{tool})}{L_0}
\]

With \emph{T} = 3: β(3) = (3 × 220) / 120 ≈ 5.5× slower than
single-shot. With \emph{T} = 1: β(1) ≈ 1.8×. The trade-off is that the
agentic system may deliver correct answers where single-shot fails, but
the price is a 2--6× latency multiplier depending on turn count.

These derived quantities give the architect concrete numbers to set
policies: cap \emph{T} at 3 unless the quality gain is statistically
significant, and allocate \textasciitilde220\,ms per turn in the latency
budget.

\section{Common Mistakes}\label{common-mistakes}

\begin{enumerate}
\def\labelenumi{\arabic{enumi}.}
\item
  \textbf{Unbounded turn limits.} Setting the maximum number of agent
  turns too high (e.g., 10 or more) leads to runaway token consumption
  and latency spikes. In production we have seen a single request
  consume \textgreater30,000 tokens and take \textgreater5\,seconds
  because the model kept looping on a poorly scoped tool.
\item
  \textbf{Neglecting token accounting per turn.} Many architects track
  only the initial input tokens and forget that each retrieval or tool
  call adds context that persists for all subsequent turns. The
  compounding effect means that turn \emph{T} sees an input base that
  grows linearly with \emph{T}.
\item
  \textbf{Overlooking tool latency variance.} A tool that typically
  responds in 30\,ms may, under cache misses or lock contention, take
  2\,s. Without a timeout and fallback, the entire agent loop blocks.
  Always instrument timeouts and provide a ``best-effort'' fallback
  path.
\item
  \textbf{Confusing `agentic' with `more prompts'.} Adding extra prompts
  without a structured loop (Plan/Execute/Observe/Decide) does not give
  the model reasoning power; it just inflates token count. The agentic
  value comes from the model's ability to reinterpret prior observations
  and decide the next action.
\item
  \textbf{Skipping guardrails.} Without pre-call validation (e.g., query
  sanitization, budget checks) and post-call verification (answer
  relevance, hallucination detection), agentic systems can produce
  plausible but incorrect or unsafe outputs. Guardrails are not optional
  --- they are the safety net that makes unbounded loops acceptable in
  production.
\end{enumerate}

\section{Architecture Consequence}\label{architecture-consequence}

Introducing an agentic layer reshapes the system architecture in four
ways:

\textbf{Inference scaling.} The base model (70B FP16, 8× H100) must now
serve variable-length contexts, and --- critically --- the KV-residency
conclusion is far more severe than a per-latency convenience because the
canonical per-token KV is \textasciitilde2.5 MB/token (\hyperref[chap:7]{Chapter 7}). A
request that triggers 4 agent turns carries \textasciitilde12,960 tokens
vs.~9,500 for single-shot (\hyperref[tab:19.1]{Table 19-1}) --- that is not a ``35\% GPU
memory'' tweak but a move from \textasciitilde23.8 GB KV per request
(9,500 tokens × 2.5 MB) to \textasciitilde32.4 GB (12,960 × 2.5 MB) ---
roughly a 36\% growth, on top of an already-large KV footprint. At 100
concurrent users with the median 2-turn profile (\textasciitilde28.1
GB/request at 11,230 tokens), the aggregate KV alone is
\textasciitilde2.7 TB --- far beyond one host's \textasciitilde436 GB KV
budget, so \textbf{agents do not just stress memory, they force a fleet}
(\hyperref[chap:17]{Chapter 17}). The architect must therefore treat context growth as a
first-order capacity driver: cap \emph{T}, cache prefixes, or quantize
KV (FP8), as these levers repeatedly outprice adding hosts. Note we must
be precise: the per-request KV is set by the model's canonical per-token
KV, not by the small per-layer figures an earlier draft used; the honest
number is tens of GB per request, which is exactly why agentic depth is
a fleet-sizing problem.

\begin{figure}
\centering
\pandocbounded{\includegraphics[keepaspectratio,alt={\hyperref[fig:19.2]{Fig 19.2} --- Input/output tokens and KV cache per request across agent turns (I₀ + T·δ + T·γ; 9,200 + 800 + 65 per turn), FP16 KV rising from 23.8 GB single-shot to 32.4 GB at 4 turns {[}DERIVED: Ch19 §3 arithmetic{]}},width=\maxwidth{\textwidth},keepaspectratio]{figures/fig-19-1902.pdf}}
\caption{\hyperref[fig:19.2]{Fig 19.2} --- Input/output tokens and KV cache per request
across agent turns (I₀ + T·δ + T·γ; 9,200 + 800 + 65 per turn), FP16 KV
rising from 23.8 GB single-shot to 32.4 GB at 4 turns {[}DERIVED: Ch19
§3 arithmetic{]}}

\label{fig:19.2}\end{figure}

\emph{\hyperref[fig:19.2]{Fig 19.2} --- Why agentic depth is a fleet-sizing problem. Both the
token count and the FP16 KV cache grow linearly with turns; four turns
take a request from \textasciitilde23.8 GB to \textasciitilde32.4 GB of
KV --- a 36\% jump on an already-large footprint, which is exactly why a
fleet (Ch17), prefix caching, or FP8 KV pay for themselves on agentic
workloads.}

\begin{figure}
\centering
\pandocbounded{\includegraphics[keepaspectratio,alt={\hyperref[fig:19.3]{Fig 19.3} --- The agentic context block grows every turn: initial prompt (9.2K tok) plus tool output + reasoning appended each turn, KV footprint rising 23.8 → 32.4 GB {[}2° DERIVED: Ch19 §3 arithmetic{]}},width=\maxwidth{\textwidth},keepaspectratio]{figures/fig-19-1903.pdf}}
\caption{\hyperref[fig:19.3]{Fig 19.3} --- The agentic context block grows every turn:
initial prompt (9.2K tok) plus tool output + reasoning appended each
turn, KV footprint rising 23.8 → 32.4 GB {[}2° DERIVED: Ch19 §3
arithmetic{]}}

\label{fig:19.3}\end{figure}

\emph{\hyperref[fig:19.3]{Fig 19.3} --- Agentic context accumulation as a memory sequence.
Each turn appends retrieved context (δ ≈ 800 tokens) plus generated
reasoning (γ ≈ 65 tokens) to a context that persists across every later
turn; the FP16 KV block therefore grows from \textasciitilde23.8 GB
(single-shot) to \textasciitilde32.4 GB (four turns) --- a 36\%
footprint rise driven purely by cumulative context, not by extra
requests. This is the mechanism behind the fleet-sizing argument in \hyperref[fig:19.2]{Fig 19.2}.}

\textbf{Task economics.} The deepest shift an agentic workload imposes
is a change in \emph{what the architect optimizes}. For conventional
inference, the unit of cost is the token: the architect optimizes token
economics --- cost per token, latency per token, KV per token. For an
agentic system, the unit of value is the \emph{task}: the architect
optimizes task economics, because a task consumes a variable,
data-dependent number of model calls, tool calls, and tokens, and may
fail. The right denominator is therefore \textbf{cost per successful
task}, not cost per token:

\begin{verbatim}
cost / successful task ≈ (model calls/task × tokens/call × $/token)
                            + (tool calls/task × $/tool)
                            + (retries/task × cost/retry)
             divided by the success rate of the task
\end{verbatim}

where tokens/call is itself inflated by the agentic amplification
\emph{α} from Section 4. Two consequences follow. First, a more
expensive model that completes tasks in one turn can be cheaper
\emph{per successful task} than a cheap model that needs three turns ---
the amplification term can dominate the base price. Second, the success
rate is as important as the unit cost: a 10\% cheaper pipeline that
fails 5\% more often can be net more expensive once retries and human
intervention are counted. This is a genuinely different objective from
\emph{cheapest tokens}, and it is why agentic workloads cannot be priced
or capacity-planned with the single-shot token arithmetic alone. An
architect should instrument tasks (success, turns, calls, retries) and
optimize cost per successful task, exactly as an architect of a
deterministic service optimizes cost per completed request.

\textbf{Tool backend provisioning.} Search indexes, SQL pools, and code
interpreters become first-class infrastructure components. Their
capacity (concurrent queries, connection pool size, query throughput)
must be dimensioned alongside the LLM GPU pool. A bottleneck in the tool
backend will manifest as increased λ\_tool, which directly inflates the
per-turn latency \emph{L} and therefore the overall request latency.

\textbf{Observability requirements.} Each turn generates a tuple of
(turn number, model input tokens, tool call, tool latency, tool output
tokens, decision). Correlating these tuples across turns for a single
request, and across requests for the fleet, requires structured logging
and tracing. Without it, debugging why a request took 2\,s and consumed
15,000 tokens is nearly impossible.

\textbf{Fallback \& degradation pathways.} Because the agentic layer
introduces data-dependent depth, the system must have graceful
degradation: if the turn limit is reached, return the best partial
answer so far, or fall back to a simpler single-shot prompt. This
preserves user experience even when the agent fails to converge.

\section{What We Still Don't Know}\label{what-we-still-dont-know}

Despite the quantified arithmetic above, several questions remain open
and would benefit from targeted research:

\begin{itemize}
\tightlist
\item
  \textbf{Turn-count distribution shift with model scale.} Do
  larger-context models (e.g., 405B FP16) exhibit different \emph{T}
  distributions? Does the ability to ``see more'' reduce the need for
  multi-step loops, or does the model reason more deeply and actually
  increase \emph{T}?
\item
  \textbf{Token amplification vs.~answer quality curve.} We have
  measured \emph{α(T)}, but we have not firmly established the marginal
  quality gain per additional turn. Is the 1.27× → 1.36× move from
  \emph{T} = 3 to \emph{T} = 4 worth the 7\% token increase? A/B tests
  across query categories would help.
\item
  \textbf{Latency tail under spike load.} Our 95th-percentile latency of
  \textasciitilde880\,ms assumes GPU inference time is constant. Under
  traffic spikes, GPU saturation can increase inference latency
  non-linearly. Empirical measurement under controlled load is needed to
  validate the linear \emph{T·L} model.
\item
  \textbf{Optimal tool design for minimal turn count.} Some tools could
  return more comprehensive results in a single call, reducing the need
  for multiple rounds. The design space of ``rich vs.~narrow'' tools and
  its impact on \emph{T} is underexplored.
\item
  \textbf{Cross-user context caching.} If many users ask related
  questions (e.g., ``What does the policy say about X?''), can we cache
  intermediate retrieval results and avoid redundant tool calls? The
  savings could be substantial but require careful invalidation logic.
\end{itemize}

\section{End-of-Chapter Mini-Case}\label{end-of-chapter-mini-case}

\textbf{Scenario.} A enterprise RQA system serves 2,000 knowledge
workers. The baseline single-shot configuration uses a 70B FP16 model on
8× H100, with 9,200 input tokens and 300 output tokens, handling 10 rps
average / 40 rps peak. The team introduces an agentic layer with a
maximum of 3 turns, each turn adding a search tool (800 tokens) and an
SQL query tool (350 tokens). The guardrail budget caps total input at
13,000 tokens.

\textbf{Measurement results.} Telemetry over 30 days shows: - 68\% of
requests terminate in 1 turn - 25\% terminate in 2 turns - 5\% terminate
in 3 turns - 2\% exceed the turn limit and fall back to single-shot

\textbf{Token and latency impact.} - Average total input tokens: 9,200 +
(0.68×800) + (0.25×(800+350)) + (0.05×(800+350+800)) ≈ 9,200 + 544 +
287.5 + 97.5 ≈ 10,129 tokens - Amplification \emph{α}: 10,129 / 9,500 ≈
1.07× - Average per-turn latency: 220\,ms → median end-to-end latency: 2
× 220\,ms ≈ 440\,ms (vs.~\textasciitilde120\,ms single-shot) -
95th-percentile latency (T ≈ 4 with fallback): \textasciitilde880\,ms

\textbf{Architectural actions.} The team sets the turn limit to 3,
instruments each turn's token and latency, and adds a context cache for
the search tool. With the cache hit rate of 30\% on recurring queries,
the effective \emph{δ} drops to \textasciitilde560 tokens, reducing
average amplification to \textasciitilde1.04× and median latency to
\textasciitilde350\,ms. The system now meets the SLA of \textless1\,s
95th-percentile latency while delivering higher-quality answers on
complex queries.

\textbf{Lesson.} The agentic layer added \textasciitilde7\% token
overhead and \textasciitilde320\,ms latency on median, but improved
answer correctness on multi-step reasoning queries by an estimated 22\%
(measured by human eval). The trade-off was acceptable, and the token
overhead was mitigated by context caching. The architect's key decisions
were: (a) capping turns at 3, (b) adding a search cache, and (c)
providing a single-shot fallback when the limit is hit.

\begin{center}\rule{0.5\linewidth}{0.5pt}\end{center}

\textbf{\hyperref[tab:19.1]{Table 19-1}} --- Token amplification and latency trade-offs for
agentic RQA pipelines

\begin{longtable}[]{@{}
  >{\raggedright\arraybackslash}p{0.1311\linewidth}
  >{\raggedright\arraybackslash}p{0.0820\linewidth}
  >{\raggedright\arraybackslash}p{0.2295\linewidth}
  >{\raggedright\arraybackslash}p{0.0820\linewidth}
  >{\raggedright\arraybackslash}p{0.2623\linewidth}
  >{\raggedright\arraybackslash}p{0.2131\linewidth}@{}}
\toprule\noalign{}
\begin{minipage}[b]{\linewidth}\raggedright
Metric
\end{minipage} & \begin{minipage}[b]{\linewidth}\raggedright
T=1
\end{minipage} & \begin{minipage}[b]{\linewidth}\raggedright
T=2 (median)
\end{minipage} & \begin{minipage}[b]{\linewidth}\raggedright
T=3
\end{minipage} & \begin{minipage}[b]{\linewidth}\raggedright
T=4 (90th pct)
\end{minipage} & \begin{minipage}[b]{\linewidth}\raggedright
Single-shot
\end{minipage} \\
\midrule\noalign{}
\endhead
\bottomrule\noalign{}
\endlastfoot
Input tokens (9,200 + T·800 + T·65) & 10,065 & 10,930 & 11,795 & 12,660
& 9,200 \\
Total tokens (incl.~300 output) & 10,365 & 11,230 & 12,095 & 12,960 &
9,500 \\
Token amplification \emph{α} & 1.09× & 1.18× & 1.27× & 1.36× & 1.00× \\
Per-turn LLM inference, 10K ctx & \textasciitilde120 ms &
\textasciitilde120 ms & \textasciitilde120 ms & \textasciitilde120 ms &
\textasciitilde120 ms \\
Tool latency (median) & \textasciitilde100 ms & \textasciitilde100 ms &
\textasciitilde100 ms & \textasciitilde100 ms & --- \\
Cumulative end-to-end latency & \textasciitilde220 ms &
\textasciitilde440 ms & \textasciitilde660 ms & \textasciitilde880 ms &
\textasciitilde120 ms \\
Latency amplification \emph{β} = T·220/120 & \textasciitilde1.8× &
\textasciitilde3.7× & \textasciitilde5.5× & \textasciitilde7.3× &
1.00× \\
KV cache @ 2.5 MB/token & \textasciitilde25.9 GB & \textasciitilde28.1
GB & \textasciitilde30.2 GB & \textasciitilde32.4 GB &
\textasciitilde23.8 GB \\

\label{tab:19.1}\end{longtable}

\emph{All token counts use δ=800, γ=65, O=300; latency uses λ\_inf≈120
ms, λ\_tool≈100 ms. These are canonical-model {[}DERIVED{]} values, not
measured production figures.}
